% Source pp. 102-110 (Lectures7-12.pdf).
\section{Renormalized Perturbation Theory}

In $\varphi^4$ theory, the vertices with positive superficial degree of
divergence $D\ge0$ (the so called \emph{primitive divergences}) are
$\Gamma^{(2)}$ and $\Gamma^{(4)}$. We have already reviewed the simplest
divergent diagram for $\Gamma^{(2)}$. The diagrams for $\Gamma^{(4)}$ have the
superficial degree of divergence $D=0$; this means these diagrams are
\emph{logarithmically divergent}. To clarify the meaning of this term, let us
consider the simplest connected diagrams for this vertex,
\[
-\tilde\Gamma^{(4)}(p_1,p_2,p_3,p_4)=
\diac{\foreach \a/\l in {135/p_1,45/p_2,315/p_3,225/p_4} {\draw[prop,densely dashed] (0,0) -- (\a:1.6);
  \node at ($(\a:1.7)+(\a+90:0.45)$) {\scriptsize$\l$};} \node[vtx] at (0,0) {};}
+\Bigl(\diac{\draw[prop] (0,0) ellipse (1.2 and 0.7);
  \foreach \a/\l in {135/p_1,225/p_4} {\draw[prop,densely dashed] (-1.2,0) -- ++(\a:1.4); \node at ($(-1.2,0)+(\a:1.6)+(\a+90:0.4)$) {\scriptsize$\l$};}
  \foreach \a/\l in {45/p_2,315/p_3} {\draw[prop,densely dashed] (1.2,0) -- ++(\a:1.4); \node at ($(1.2,0)+(\a:1.6)+(\a-90:0.4)$) {\scriptsize$\l$};}
  \draw[arr] (-0.4,1.0) -- (0.4,1.0); \node[above] at (0,1.05) {\scriptsize$p_1+p_2-k$};
  \draw[arr] (-0.4,-1.0) -- (0.4,-1.0); \node[below] at (0,-1.05) {\scriptsize$k$};
  \node[vtx] at (-1.2,0) {}; \node[vtx] at (1.2,0) {};}
+\text{permutations}\Bigr)+\cdots
\]
The explicitly shown loop diagram represents the momentum integral
\[
\frac{(-\lambda)^2}2\int\frac{d^4k}{(2\pi)^4}\,\frac1{(k^2+m^2)\bigl((p_1+p_2-k)^2+m^2\bigr)}\,.
\]
When $|k|\gg|p_1+p_2|$ and $|k|\gg m$ this reduces to
\[
\frac{\lambda^2}{2\,(2\pi)^4}\int\frac{d^4k}{k^2\,k^2}\,.
\]
Note that $d^4k=d\Omega_4\,|k|^3\,d|k|$, where $d\Omega_4$ is the surface
element of a unit 3-sphere $S^3$, $\int_{S^3}d\Omega_4=2\pi^2$. We have
\[
\frac{\lambda^2\,2\pi^2}{2\,(2\pi)^4}\int\frac{|k|^3\,d|k|}{|k|^4}=\frac{\lambda^2}{16\pi^2}\int\frac{d|k|}{|k|}\,.
\]
The integral diverges at large $|k|$ (one shouldn't worry about its divergence
at $|k|\to0$; the above simplified form is not valid at small $|k|$). If one
introduces a cutoff, say replaces
\[
\frac1{k^2+m^2}\ \to\ \frac1{k^2+m^2}\,\Phi(k^2/\Lambda^2)
\]
with some $\Phi(x)$ such that $\Phi(x)=1$ for $x\ll1$, but decays fast at
$x\gg1$, the integral evaluates as
\[
\frac{\lambda^2}{32\pi^2}\Bigl(\log\frac{\Lambda^2}{m^2}+f(p_{12})\Bigr)\,,
\]
where $f(p_{12})$ (I use $p_{12}=p_1+p_2$) has finite limit as
$\Lambda\to\infty$. Including the ``permutations'' terms, we then have
\[
\tilde\Gamma^{(4)}=\lambda-\frac{\lambda^2}{32\pi^2}\Bigl(3\log\frac{\Lambda^2}{m^2}+f(p_{12})+f(p_{13})+f(p_{14})\Bigr)\,.
\]

One can observe that the divergent term carries no dependence on the external
momenta, and in this respect comes on the same footing as the leading
contribution $\lambda$. Therefore at least at the order $\lambda^2$ the
$\Lambda$-dependent term can be absorbed into some redefinition of the coupling
constant $\lambda$. This can be formulated more precisely as follows. Just as
in the case of the mass renormalization discussed above, we should start with
the action
\[
\mathcal A=\int\Bigl[\frac12\,(\partial\varphi)^2+\frac{m_0^2}2\,\varphi^2+\frac{\lambda_0}{4!}\,\varphi^4\Bigr]\,,
\]
with some ``bare'' coupling constant $\lambda_0$ (as well as the ``bare'' mass
parameter $m_0^2$). Then it is possible to give $\lambda_0$ certain dependence
on the cutoff momentum $\Lambda$,
\[
\lambda_0=\lambda_0(\Lambda)\,,
\]
such that the limit $\Lambda\to\infty$ of $\tilde\Gamma^{(4)}$ in this order of
the perturbation theory is finite. Namely, let us set
\[
\lambda_0(\Lambda)=\lambda+\frac{3\lambda^2}{32\pi^2}\Bigl(\log\frac{\Lambda^2}{m^2}+C\Bigr)+O(\lambda^3)\,,
\]
where new parameter $\lambda$ is called the \emph{renormalized} coupling
constant, and $C$ is arbitrary number. We then have
\begin{align*}
\tilde\Gamma^{(4)}&=\lambda_0-\frac{\lambda_0^2}{32\pi^2}\Bigl(3\log\frac{\Lambda^2}{m^2}+\tilde f(p)\Bigr)+O(\lambda_0^3)\\
&=\lambda+\frac{3\lambda^2}{32\pi^2}\Bigl(\log\frac{\Lambda^2}{m^2}+C\Bigr)-\frac{\lambda^2}{32\pi^2}\Bigl(3\log\frac{\Lambda^2}{m^2}+\tilde f(p)\Bigr)+O(\lambda^3)\,.
\end{align*}
(Here $\tilde f(p)=f(p_{12})+f(p_{13})+f(p_{14})$.) The $\Lambda$-dependent term
cancels out. Assuming that $\lambda$ is fixed in the limit $\Lambda\to\infty$,
the above proper vertex has finite limit when the cutoff is removed.

Later we shall formulate this renormalization prescription with full
precision, and in particular explain the meaning of this constant $C$. But for
now we still have one more divergence to deal with.

Consider the diagram
\[
\diac{\draw[prop,densely dashed] (-2.4,0) -- (-1.2,0) (1.2,0) -- (2.4,0);
  \draw[prop] (-1.2,0) -- (1.2,0); \draw[prop] (0,0) ellipse (1.2 and 0.75);
  \node[above] at (-2.0,0) {\small$p$};
  \draw[arr] (-0.4,1.05) -- (0.4,1.05); \node[above] at (0,1.1) {\small$p-k_1-k_2$};
  \draw[arr] (-0.4,0.25) -- (0.4,0.25); \node[below] at (0,0.05) {\scriptsize$k_2$};
  \draw[arr] (-0.4,-1.05) -- (0.4,-1.05); \node[below] at (0,-1.1) {\small$k_1$};
  \node[vtx] at (-1.2,0) {}; \node[vtx] at (1.2,0) {};}\,,
\]
contributing to $-\tilde\Sigma(p)$ in the order $\lambda^2$. I will denote this
contribution as $-\tilde\Sigma_2(p)$. According to the general counting above,
this diagram has superficial degree of divergence $D=2$, i.e.\ with the cutoff
introduced it behaves as $\Lambda^2$ as $\Lambda\to\infty$. We can see this
writing the contribution explicitly,
\[
-\tilde\Sigma_2(p)=\frac{(-\lambda)^2}{3!}\int\frac{d^4k_1}{(2\pi)^4}\,\frac{d^4k_2}{(2\pi)^4}\,
\frac1{(k_1^2+m^2)(k_2^2+m^2)\bigl((p-k_1-k_2)^2+m^2\bigr)}
\]
(note the symmetry factor $1/3!$). For $k_1,k_2\gg p,m$ the integrand
simplifies as
\[
\sim\int\frac{d^4k_1\,d^4k_2}{k_1^2\,k_2^2\,(k_1+k_2)^2}\sim\Lambda^2\,.
\]
Let us write
\[
\tilde\Sigma_2(p)=\tilde\Sigma_2(0)+\bigl(\tilde\Sigma_2(p)-\tilde\Sigma_2(0)\bigr)\,.
\]
Here
\[
\tilde\Sigma_2(0)=-\frac{\lambda^2}{3!}\int\frac{d^4k_1\,d^4k_2}{(2\pi)^8}\,\frac1{(k_1^2+m^2)(k_2^2+m^2)\bigl((k_1+k_2)^2+m^2\bigr)}\,.
\]
This divergent term does not depend on $p$; in this respect it is similar to
the contribution of the diagram
\[
\dia{\draw[prop,densely dashed] (-2,0) -- (2,0); \draw[prop] (0,0.6) circle (0.6); \node[vtx] at (0,0) {};}
\]
Clearly, it also can be absorbed into the mass renormalization, by suitable
modification of the counterterm
\[
\frac{\delta m^2(\Lambda)}2\,\varphi^2\,.
\]
The difference
\begin{multline*}
\tilde\Sigma_2(p)-\tilde\Sigma_2(0)\\
=-\frac{\lambda^2}{3!}\int\frac{d^4k_1\,d^4k_2}{(2\pi)^8}\,\frac1{(k_1^2+m^2)(k_2^2+m^2)}\,
\frac{2(pk_1)+2(pk_2)-p^2}{\bigl((k_1+k_2)^2+m^2\bigr)\bigl((k_1+k_2-p)^2+m^2\bigr)}
\end{multline*}
still diverges (logarithmically). If the cutoff is introduced, this
contribution is written as
\[
\sim p^2\log\frac{\Lambda^2}{m^2}+\text{finite}\,.
\]
This divergent term cannot be absorbed by the mass renormalization, since it
depends on the momentum $p$. Note however that this dependence has simple form
$\sim p^2$.

Recall that the proper vertex $\Gamma^{(2)}$ has the form
\[
\tilde\Gamma^{(2)}(p)=p^2+m^2+\text{loop corrections}\,,
\]
where the $p^2$ term originates from the term
\[
\frac12\,(\partial\varphi)^2
\]
term in the original action. This suggests that the above divergent term
$\sim p^2$ can be understood in terms of the ``field renormalization''.
Originally, the field $\varphi$ appeared in the theory as the integration
variable in the functional integral. Let us rename this integration variable as
$\varphi_0$, and call it the ``bare field''. The new divergent term shows that
the correlation functions of this bare field retains the dependence on the
cutoff parameter $\Lambda$, even after the mass and coupling renormalizations
are carried out. The way this dependence appears (through the $p^2$ term)
suggests that one can define renormalized field,
\[
\varphi(x)=Z^{-1/2}(\Lambda)\,\varphi_0(x)\,,
\]
with certain cutoff-dependent renormalization constant $Z(\Lambda)$, such that
the correlation functions
\[
\bigl\langle\varphi(x_1)\cdots\varphi(x_n)\bigr\rangle
\]
admit finite $\Lambda\to\infty$ limit.

We are ready to formulate the \emph{renormalization program}.

\subsection*{Renormalization program (\texorpdfstring{$\varphi^4$}{phi**4} theory)}

1. We start with the action
\[
\mathcal A=\int d^4x\,\Bigl[\frac12\,(\partial\varphi_0)^2+\frac{m_0^2}2\,\varphi_0^2+\frac{\lambda_0}{4!}\,\varphi_0^4\Bigr]\,,
\]
which involves the bare field $\varphi_0$, the bare mass $m_0$, and the bare
coupling constant $\lambda_0$.

2. We introduce some cutoff, with the cutoff momentum $\Lambda$, to make all
diagram contributions finite but cutoff-dependent (this can be done in many
ways, as we shall discuss later).

3. The bare parameters $m_0^2$, $\lambda_0$, and the field renormalization
constant $Z$ should be given certain dependence on the cutoff momentum
$\Lambda$,
\[
m_0^2(\Lambda)\,,\qquad\lambda_0(\Lambda)\,,\qquad Z(\Lambda)\,,
\]
such that the correlation functions of the renormalized field
$\varphi(x)=Z^{-1/2}\varphi_0(x)$ will have finite $\Lambda\to\infty$ limit.

In practice, one uses the perturbation theory to evaluate the correlation
functions. The above program can be further reformulated as the

\subsection*{Renormalized perturbation theory}

Let us rewrite the action as
\[
\mathcal A=\int d^4x\,\Bigl[\underbrace{\frac12\,(\partial\varphi)^2+\frac{m^2}2\,\varphi^2}_{\text{free part}}
+\underbrace{\frac\lambda{4!}\,\varphi^4}_{\text{interaction part}}
+\underbrace{\frac{\delta Z}2\,(\partial\varphi)^2+\frac{\delta m^2}2\,\varphi^2+\frac{\delta\lambda}{4!}\,\varphi^4}_{\text{counterterms}}\Bigr]\,,
\]
where $m$ is the actual mass and $\lambda$ is suitably defined renormalized
coupling constant (to be discussed later), and $\varphi$ is the renormalized
field. The last three terms are the \emph{counterterms}. The identity with the
original action implies
\[
1+\delta Z=Z\,,\qquad m^2+\delta m^2=Z\,m_0^2\,,\qquad\lambda+\delta\lambda=Z^2\,\lambda_0\,.
\]
In perspective, the renormalized perturbation theory is an expansion in the
renormalized coupling constant $\lambda$. We should assume that the
counterterm coefficients $\delta Z$, $\delta m^2$, and $\delta\lambda$
themselves depend perturbatively (i.e.\ as the power series) on $\lambda$,
\[
\delta Z=z_2\,\lambda^2+z_3\,\lambda^3+\cdots\,,\qquad
\delta m^2=b_1\,\lambda+b_2\,\lambda^2+\cdots\,,\qquad
\delta\lambda=a_2\,\lambda^2+a_3\,\lambda^3+\cdots\,,
\]
with $\Lambda$-dependent coefficients. We must show that these coefficients can
be adjusted in such a way that the correlation functions of the renormalized
field $\varphi$ are $\Lambda$-independent order by order in $\lambda$.

Let me stress that this renormalization program is based on our observation
that the $\varphi^4$ perturbation theory contains only finitely many basic
types of divergent diagrams (the ``primitive divergences''). Namely, the
primitively divergent diagrams are in $\Gamma^{(2)}$ and $\Gamma^{(4)}$. The
diagrams for higher vertices $\Gamma^{(n)}$ can be divergent only through
subdiagrams of the $\Gamma^{(2)}$ and $\Gamma^{(4)}$ types. The idea of the
renormalization program is that the three primitive divergences discussed above
can be absorbed in three renormalizations
\[
m_0^2(\Lambda)=Z^{-1}(\Lambda)\,\bigl(m^2+\delta m^2(\Lambda)\bigr)\,;\qquad
\lambda_0(\Lambda)=Z^{-2}(\Lambda)\,\bigl(\lambda+\delta\lambda(\Lambda)\bigr)\,;\qquad
\varphi_0=Z^{\frac12}(\Lambda)\,\varphi\,,
\]
order by order in the renormalized coupling constant $\lambda$.

Before we turn to execution of this program in $\varphi^4$ theory let us briefly
discuss divergences and renormalizations in a scalar field theory whose
interaction term is more general polynomial in $\varphi$.

\subsection*{Divergences in Scalar Theories}

It is instructive to consider such field theory in the space of $d$
dimensions, with generic $d$. We will consider general scalar theories with the
actions of the form
\begin{equation}\label{eq:genscalar}
\mathcal A=\int d^dx\,\Bigl[\frac12\,(\partial\varphi_0)^2+\frac{m_0^2}2\,\varphi_0^2+\sum_{k=3}^N\frac{\lambda_{0\,k}}{k!}\,\varphi_0^k\Bigr]\,.
\end{equation}
The Feynman rules remain the same as in 4-dimensional theory, except for the
momentum integrations go over $d$-dimensional momentum space,
\[
\frac{d^4k}{(2\pi)^4}\ \to\ \frac{d^dk}{(2\pi)^d}\,,
\]
and of course the diagrams contain $k$-leg vertices associated with the
couplings $\lambda_{0\,k}$.

As before, we can consider a generic diagram contributing to
$\tilde\Gamma^{(n)}$, containing $P$ propagators and $I$ $d$-momentum
integrations. Because now the momenta are $d$-dimensional, the superficial
degree of divergence of this diagram is
\[
D=d\cdot I-2\cdot P\,.
\]
Assume that the diagram contains $V_k$ $k$-leg vertices. Analysis similar to
that we have made for $\varphi^4$ theory reveals two identities
\[
2P+n=\sum_kk\cdot V_k\,,\qquad I=P-\sum_kV_k+1\,.
\]
Combining these equations one finds
\[
D=\sum_k\Bigl(\frac{d-2}2\,k-d\Bigr)V_k-\frac{d-2}2\,n+d\,.
\]
This formula admits simple interpretation in terms of dimensional counting.
Note that in our units $c=\hbar=1$ there is only one independent unit, which we
take to be mass unit. I will denote by $[X]$ the mass dimension of a quantity
$X$, for example
\[
[\textit{mass}]=1\,,\qquad[\textit{length}]=-1\,.
\]
The action is dimensionless and therefore it follows from its form that
\begin{equation}\label{eq:engdim}
[\varphi_0]=\frac{d-2}2\,,\qquad[m_0^2]=2\,,\qquad[\lambda_{0\,k}]=d-\frac{d-2}2\,k\,.
\end{equation}
I want to note in passing that this simple dimensional analysis is applied to
the \emph{bare} quantities. We will see later that because of renormalization
constant $Z(\Lambda)$ the renormalized field $\varphi$ can have different
dimension. This is why the above dimensions in \eqref{eq:engdim} are often
referred to as canonical (or ``engineering'') dimensions.

It is easy to check that
\[
\bigl[\tilde\Gamma^{(n)}\bigr]=d-\frac{d-2}2\,n\,.
\]
Therefore the equation for the superficial degree can be written as
\begin{equation}\label{eq:Ddim}
D=\bigl[\tilde\Gamma^{(n)}\bigr]-\sum_k\bigl[\lambda_{0\,k}\bigr]\cdot V_k\,.
\end{equation}
Now the meaning of this equation must be obvious. For $\Lambda\gg|p_i|$ the
dominating contribution of the diagram with $V_k$ vertices $\lambda_k$ is
\begin{equation}\label{eq:GammaLambda}
\tilde\Gamma^{(n)}\sim\Bigl[\prod_k\lambda_{0\,k}^{V_k}\Bigr]\,\Lambda^D\,,
\end{equation}
and our equation simply describes the balance of dimensions.

It is clear from the Eq.~\eqref{eq:Ddim} that the mass dimensions of the
coupling constants $\lambda_{0\,k}$ must play the key role in the analysis of
the perturbative divergences. Suppose some coupling $\lambda_{0\,k}$ in the
action have strictly negative mass dimension, $[\lambda_{0\,k}]<0$. Then there
are divergent contributions to $\tilde\Gamma^{(n)}$ with any $n$ which come
from the diagrams with sufficiently large $V_k$. In other words such theory has
infinitely many primitive divergences which of course cannot be absorbed by any
finite number of counterterms. Quantum field theories of this type are called
(perturbatively) \emph{non-renormalizable}. Overall consistency of
non-renormalizable theories is questionable (but not impossible). From pure
pragmatic point of view, the necessity to introduce infinitely many
counterterms brings in also infinitely many free parameters, and predictive
power of such theories is limited.

If mass dimensions of all coupling constants in the action are non-negative,
the equation \eqref{eq:Ddim} shows that there is only finite number of
primitively divergent proper vertices (this is true if $d>2$, the case $d=2$ is
exceptional and must be studied separately); these divergences then can be
absorbed by finitely many renormalization constants. For this reason such
theories are called (perturbatively) \emph{renormalizable}. Moreover, if all
$\lambda_{0\,k}$ have strictly positive mass dimensions there is only finite
number of divergent diagrams. The theories of this last kind are often referred
to as \emph{super-renormalizable}, in which case the term ``renormalizable'' is
reserved for the theories with infinitely many divergent diagrams but finitely
many primitive divergences; in view of \eqref{eq:Ddim} this requires that at
least one of the coupling constants is dimensionless. We will see later when
studying renormalization group that super-renormalizable scalar theories (with
the polynomial $V(\varphi_0)$ bounded from below) are consistent field
theories. Overall consistency of renormalizable theories requires more subtle
analysis, but at least they make sense perturbatively.

One observes from \eqref{eq:engdim} that for $d>6$ there are no (interacting)
renormalizable scalar theories. For $6\ge d>2$ renormalizable field theories
are
\[
\begin{array}{ll}
\varphi^3 & \text{in}\quad d=6\quad\text{dimensions,}\\
\varphi^3,\ \varphi^4 & \text{in}\quad d=4\quad\text{dimensions,}\\
\varphi^3,\ \varphi^4,\ \varphi^5,\ \varphi^6 & \text{in}\quad d=3\quad\text{dimensions.}
\end{array}
\]
It is not difficult to show that any polynomial in $\varphi$ is
super-renormalizable in $d=2$ dimensions.

I want to stress here that nonexistence of perturbatively renormalizable field
theories in high dimensionalities does not imply that consistent field
theories are limited to low space-time dimensions; there may exist perfectly
consistent field theories which are just ``too far'' from free field theory to
admit meaningful perturbative interpretation. I hope to discuss this
possibility in greater detail later.
